\subsection{Shape Basis}\label{sec:shape-basis}

The complete particular solution \eqref{eq:iesan-general-solution} is represented by a polynomial in the axial coordinate:
\begin{equation*}
\hat{\bm{u}}_{\rm P}(x,\bm{r}) = \sum_{m =0}^3 x^m \mathbf{S}_m(\bm{r})
\end{equation*}
where $\mathbf{S}_m(\bm{r})$ collects all sectional patterns that multiply $x^m$. 
The kinematic fields $\bm{u}(x)$ and $\bm{\theta}(x)$ generate a subspace of patterns representing plane-section deformations:
\begin{equation*}
\PlaneMotions := \big\{\bm{v}(x) + \bm{\omega}(x) \times \bm{r} \mid \bm{v}, \bm{\omega} \text{ arbitrary smooth functions of } x\big\}.
\end{equation*}
At each value of $\reflenx$, this subspace is spanned pointwise by the plane-section basis:
\begin{equation}
\PlaneMotions_x(\Section) = \operatorname{span}\{\FrameBasis,\,\FrameBasis_y,\, \FrameBasis_z,\, \FrameBasis \times \bm{r}, 
\FrameBasis\,\ry,\,
\FrameBasis\,\rz\,
%\FrameBasis \otimes \bm{r}
\}
= \operatorname{span} \begin{bmatrix}
    \mathbf{A}_{u} & \mathbf{A}_{\theta} 
\end{bmatrix}\bm{\eta}(x)
\label{eq:rigid-span}
\end{equation}
with \(y = \FrameBasis_{y} \cdot\bm{r}\) and \(z = \FrameBasis_{z}\cdot\bm{r}\). 
% \begin{Aside}
and
\begin{subequations}
    \begin{align}
        % \bm{\Psi}_{\varepsilon} &= \FrameBasis \\
        \mathbf{A}_{u}    &= \mathbf{1} \\ % \oneS
        \mathbf{A}_{\theta} & = -\bm{r}^{\times}
        % \bm{\Psi}_\vartheta     &= \FrameBasis\times \bm{r} \\
        % \bm{\Psi}_{\kappa}      &= \FrameBasis\otimes\bm{r}
    \end{align}
    \label{eq:rigid-basis}
\end{subequations}
% \end{Aside}
These functions collectively represent all plane-section motions of the form \(\bm{u} + \bm{\theta}\times\bm{r}\). 
% Any displacement pattern in the span \eqref{eq:rigid-span} can be absorbed into the  fields $(\bm{u}, \bm{\theta})$ and need not be counted as warping. 
% However, we may choose to distribute these components between the beam fields \(\bm{u},\bm{\theta}\) and the warping fields \(\bm{\alpha}\), creating a redundant representation. 
% This redundancy provides flexibility in defining strain measures and achieving desired properties in the resulting beam theory.  

% \subsubsection{Warping Basis}

The warping shapes \(\bm{\Phi}\) must contain at least the $x$-independent vector fields on the section that remain after removing all patterns that can be generated by rigid sectional motion \eqref{eq:rigid-span}.
% We term this quantity the \emph{shape rank} of the representation.
% It is worth noting that the shape rank counts distinct spatial patterns across the section, which differs from counting the number of independent scalar amplitude functions $\alpha_j(x)$ needed to drive these patterns. 
% Multiple warping shapes may share the same axial variation when their amplitudes can be expressed in terms of the one-dimensional strain measures $(\epsilon, \bm{\gamma}, \bm{\kappa})$. 

Let $\mathbf{Q}(x,\bm{r})$ denote the residual obtained by projecting the Saint-Venant field onto the quotient space $\mathcal{V}_r / \mathscr{R}(\Section)$, where $\mathcal{V}_r$ is the space of vector fields on the section. 
% This residual is a function of both $x$ and $\bm{r}$, and may be viewed as an element of the tensor product $\operatorname{span}\{1, x, x^2, x^3\} \otimes (\mathcal{V}_r / \mathscr{R}(\Section))$. 
The minimum number of warping shapes is the \emph{separation rank} of this residual, defined as the smallest integer $N$ such that
\begin{equation*}
\mathbf{Q}(x,\bm{r}) = \sum_{j=1}^{N} \alpha_j(x) \bm{\Phi}_j(\bm{r})
\end{equation*}
with each $\bm{\Phi}_j$ independent of $x$. 
% \begin{proposition}[Shape Rank]
% \label{prop:shape-rank}
% For a generic cross-section with no symmetry assumptions and arbitrary centroidal location $\CentroidLocation$, the Saint-Venant family can be represented minimally with $\ShapeRank = 8$ independent warping shapes of $\bm{r}$ after removing rigid modes.
% \end{proposition}
After removing all patterns in $\mathscr{R}(\Section)$ from the Saint-Venant solution, the residual is spanned by the following linearly independent patterns:
\begin{Old}
\begin{subequations}
\begin{alignat}{8}
&\text{Transverse:}\notag\\
&&&\bm{W}_{(\varepsilon)}:=-\,\nu\,\bm{r}&&\text{ Poisson extension }& \text{1 amplitude }&&\alpha_{\varepsilon}\propto a_\varepsilon \sim \reflenx^0\\
&&&\bm{W}_{(a)}:=-\,\nu\,\operatorname{dev}\bm{r}\otimes\bm{r}&&\text{ Poisson bending }&\text{2 amplitudes }&&\bm{\alpha}_a\propto\IesanLoadM\sim \reflenx^0 \label{eq:def-Wa}\\
&&&\PoissonFlexure:=-\,\nu\big(\operatorname{dev}\bm{r}\otimes\bm{r}-\bm{r}\otimes\CentroidLocation\big)&&\text{ Poisson flexure }&\text{2 amplitudes }&&\bm{\alpha}_{b}\propto\bm{b}_{\Section}\sim \reflenx^1
\label{eq:def-Wb}
\\
%
&\text{Longitudinal:} \notag\\
&&&\FrameBasis\otimes\bm{\WarpShearIesan}&&\text{ Flexure }&\text{2 amplitudes }&&\bm{\alpha}_{\WarpShearIesan}\propto\bm{b}_{\Section}\sim \reflenx^0 \\
&&&\FrameBasis\,\WarpTwistIesan &&\text{ Torsion }&\text{1 amplitude }&&\alpha_{\WarpTwistIesan} \propto a_{\vartheta}+c_{\vartheta}\sim \reflenx^0 
\end{alignat}
\label{eq:def-poisson-warping}
\end{subequations}
\end{Old}
\begin{subequations}
\begin{alignat}{3}
&\text{Poisson extension:} \quad && \bm{W}_{(\varepsilon)} := -\nu\bm{r}, \quad && \alpha_{\varepsilon} \propto a_\varepsilon, \\
&\text{Poisson bending:} \quad && \bm{W}_{(a)} := -\nu\operatorname{dev}\bm{r}\otimes\bm{r}, \quad && \bm{\alpha}_a \propto \IesanLoadM, \\
&\text{Poisson flexure:} \quad && \PoissonFlexure := -\nu(\operatorname{dev}\bm{r}\otimes\bm{r} - \bm{r}\otimes\CentroidLocation), \quad && \bm{\alpha}_b \propto \bm{b}_{\Section}, \\
&\text{Shear warping:} \quad && \FrameBasis \otimes \bm{\WarpShearIesan}, \quad && \bm{\alpha}_{\WarpShearIesan} \propto \bm{b}_{\Section}, \\
&\text{Torsion warping:} \quad && \FrameBasis \WarpTwistIesan, \quad && \alpha_{\WarpTwistIesan} \propto a_{\vartheta} + c_{\vartheta}.
\end{alignat}
% \label{eq:def-poisson-warping}
\label{eq:def-eight-shapes}
\end{subequations}
where recall \(\nu\) is Poisson's ratio and the centroid location \(\CentroidLocation\) is defined by \eqref{eq:def-centroid}.

Linear independence of these patterns in the quotient space $\mathcal{V}_r / \PlaneMotions$ follows from the fact that no nontrivial linear combination can be written as a rigid mode. 
This is verified by examining the harmonic and Poisson problems \eqref{eq:bvp-warp-twist-iesan} and \eqref{eq:bvp-warp-shear-iesan} satisfied by $\WarpTwistIesan$ and $\bm{\WarpShearIesan}$, together with the polynomial structure of the Poisson modes, which ensures that no combination of these  patterns lies in the span of the section-rigid basis functions \eqref{eq:rigid-span}.
% \end{proof}

% This result establishes that any beam model endeavoring to exactly represent the complete Saint-Venant solution class must employ at least eight warping shapes beyond those generated by rigid sectional motion. 

